% Gerado a partir de Algoritmos_Resolucao_Sinais_Sistemas.md
\chapter{Guia Algorítmico, Macetes e Exemplos Clássicos: Sinais e Sistemas}

Este guia apresenta passos estruturados, heurísticas visuais e exemplos classificados em \textbf{Fácil}, \textbf{Médio} e \textbf{Nível de Prova}, com base nos livros clássicos e nas suas \textbf{próprias listas e provas (Prof.~Edgar Silva Júnior)}.

\emph{As resoluções foram abertas linha a linha, simulando a escrita à mão no papel.}

\section{PARTE 1: PROPRIEDADES DOS SINAIS}\label{parte-1-propriedades-dos-sinais}

\subsection{1. Sinal Par, Ímpar ou Nenhum dos dois}\label{sinal-par-uxedmpar-ou-nenhum-dos-dois}

Um sinal é par se apresenta simetria em relação ao eixo vertical, e ímpar se apresenta anti-simetria.

\textbf{Algoritmo de Resolução:}

\begin{enumerate}
\def\labelenumi{\arabic{enumi}.}
\tightlist
\item
  \textbf{Passo 1:} Escreva a expressão original do sinal \(x(t)\) (ou \(x[n]\)).
\item
  \textbf{Passo 2:} Crie um novo sinal refletido no tempo: substitua todos os \(t\) (ou \(n\)) da equação por \(-t\) (ou \(-n\)). Chame esse resultado de \(x(-t)\).
\item
  \textbf{Passo 3:} Compare \(x(-t)\) com a expressão original \(x(t)\). Se \(x(-t) == x(t)\) para todos os valores, então o sinal é \textbf{PAR}. Fim.
\item
  \textbf{Passo 4:} Se forem diferentes, pegue a expressão original \(x(t)\) e multiplique toda ela por \(-1\), obtendo \(-x(t)\).
\item
  \textbf{Passo 5:} Compare \(x(-t)\) com \(-x(t)\). Se \(x(-t) == -x(t)\), então o sinal é \textbf{ÍMPAR}. Se também for diferente, não é nem par nem ímpar.
\end{enumerate}

\begin{quote}
\textbf{Dica de Ouro / Macete Visual:}

\begin{itemize}
\tightlist
\item
  Par \(\times\) Par = Par
\item
  Ímpar \(\times\) Ímpar = Par
\item
  Par \(\times\) Ímpar = Ímpar
\end{itemize}
\end{quote}

\textbf{Exemplos Práticos:}

\begin{itemize}
\item
  \textbf{Fácil (Trivial):} \(x(t) = \cos(t)\).\\
  Substituindo o tempo:
  \[x(-t) = \cos(-t)\]
  Sendo o cosseno uma função par:
  \[x(-t) = \cos(t) = x(t)\]
  Logo, é \textbf{Par}.
\item
  \textbf{Médio:} \(x(t) = t \cdot \sin(t)\).\\
  Calculando a reflexão no tempo:
  \[x(-t) = (-t) \cdot \sin(-t)\]
  Sabendo que o seno é ímpar, ou seja, \(\sin(-t) = -\sin(t)\):
  \[x(-t) = (-t) \cdot (-\sin(t))\]
  \[x(-t) = t \cdot \sin(t)\]
  \[x(-t) = x(t)\]
  Logo, é \textbf{Par}.
\item
  \textbf{Difícil:} O Degrau Unitário \(u(t)\).\\
  A função original \(u(t)\) não é nem par nem ímpar. O \emph{Oppenheim} pede a decomposição.\\
  Parte Par:
  \[x_p(t) = \frac{u(t) + u(-t)}{2} = \frac{1}{2}\]
  Parte Ímpar:
  \[x_i(t) = \frac{u(t) - u(-t)}{2} = \frac{1}{2}\text{sgn}(t)\]
\end{itemize}

\subsection{2. Sinal de Energia ou de Potência}\label{sinal-de-energia-ou-de-potuxeancia}

\textbf{Algoritmo de Resolução:}

\begin{enumerate}
\def\labelenumi{\arabic{enumi}.}
\tightlist
\item
  \textbf{Passo 1:} Calcule a energia total do sinal \(E_\infty = \lim_{T \to \infty} \int_{-T}^{T} |x(t)|^2 dt\).
\item
  \textbf{Passo 2:} Se a integral convergir para um número finito (\(0 < E_\infty < \infty\)), então é um \textbf{Sinal de Energia}. A potência média será zero. Fim.
\item
  \textbf{Passo 3:} Se \(E_\infty = \infty\), calcule a potência média: \(P_\infty = \lim_{T \to \infty} \frac{1}{2T} \int_{-T}^{T} |x(t)|^2 dt\).
\item
  \textbf{Passo 4:} Se o limite convergir para um número finito (\(0 < P_\infty < \infty\)), é um \textbf{Sinal de Potência}.
\end{enumerate}

\begin{quote}
\textbf{Dica de Ouro / Macete Visual:}\\
Sinais que ``morrem'' em infinito são de Energia (ex: exponenciais decrescentes). Sinais que duram para sempre oscilando ou constantes são de Potência (ex: degraus e senoides).
\end{quote}

\textbf{Exemplos Práticos:}

\begin{itemize}
\item
  \textbf{Fácil (Trivial):} Exponencial decrescente \(x(t) = e^{-at}u(t)\), com \(a > 0\).\\
  Energia:
  \[E = \int_{0}^{\infty} (e^{-at})^2 dt\]
  \[E = \left[ \frac{e^{-2at}}{-2a} \right]_{0}^{\infty}\]
  \[E = \frac{1}{2a}\]
  Finito, logo é \textbf{Energia}.
\item
  \textbf{Médio:} Senóide genérica \(x(t) = A\cos(\omega_0 t)\).\\
  A integral de energia vai dar \(\infty\). Indo pra Potência num período \(T\):
  \[P = \frac{1}{T} \int_{-T/2}^{T/2} A^2 \cos^2(\omega_0 t) dt = \frac{A^2}{2}\]
  Sinal de \textbf{Potência}.
\item
  \textbf{Nível Prova/Lista (Difícil):} Energia da integral de impulsos.\\
  \emph{(Origem: Questão 3 da sua 1ª Avaliação / Q. 21 da Lista)}\\
  Seja \(x(t) = \delta(t + 2) - \delta(t - 2)\). Calcule a Energia \(E_\infty\) de \(y(t) = \int_{-\infty}^{t} x(\tau) d\tau\).\\
  Passo 1 (Resolver a integral analiticamente, lembrando que integral do impulso é o degrau):
  \[y(t) = \int_{-\infty}^{t} [\delta(\tau + 2) - \delta(\tau - 2)] d\tau\]
  \[y(t) = u(t + 2) - u(t - 2)\]
  Passo 2 (Interpretação): Isso forma um pulso retangular de amplitude 1 que existe apenas de \(t = -2\) até \(t = 2\).\\
  Passo 3 (Calcular Energia):
  \[E_\infty = \int_{-2}^{2} (1)^2 dt\]
  \[E_\infty = [t]_{-2}^{2}\]
  \[E_\infty = 2 - (-2) = 4 \text{ J}\]
  Como é finito (\(E=4\)), é \textbf{Sinal de Energia}.
\end{itemize}

\section{PARTE 2: PROPRIEDADES DOS SISTEMAS}\label{parte-2-propriedades-dos-sistemas}

Um sistema mapeia uma entrada \(x(t)\) para uma saída \(y(t)\).

\subsection{3. Com Memória ou Sem Memória (Estático)}\label{com-memuxf3ria-ou-sem-memuxf3ria-estuxe1tico}

Um sistema sem memória só depende do instante presente.

\textbf{Algoritmo de Resolução:}

\begin{enumerate}
\def\labelenumi{\arabic{enumi}.}
\tightlist
\item
  \textbf{Passo 1:} Avalie a equação de saída \(y(t)\) para um instante de tempo arbitrário \(t = t_0\).
\item
  \textbf{Passo 2:} Identifique os argumentos de tempo de todas as entradas \(x(\cdot)\) que compõem a saída.
\item
  \textbf{Passo 3:} Se você precisa de \(x(t_0 - 1)\) (passado) ou \(x(2t_0)\) ou \(x(t_0+1)\) (futuro), o sistema precisa armazenar valores. Logo, tem \textbf{MEMÓRIA} (dinâmico).
\item
  \textbf{Passo 4:} Se e somente se aparecer \textbf{apenas} \(x(t_0)\), então o sistema é \textbf{SEM MEMÓRIA} (estático).
\end{enumerate}

\begin{quote}
\textbf{Dica de Ouro / Macete Visual:}\\
Integrais, derivadas e modificadores de tempo dentro do argumento do x (como \(t-1\), \(2t\), \(-t\)) sempre exigem memória. Para ser sem memória, a relação precisa ser puramente algébrica instantânea.
\end{quote}

\textbf{Exemplos Práticos:}

\begin{itemize}
\item
  \textbf{Fácil (Trivial):} \(y(t) = 5 \cdot x(t)\).\\
  Para saber a saída em \(t=2\), só usa \(x(2)\). \textbf{Sem Memória}.
\item
  \textbf{Médio:} \(y(t) = x(t-2)\).\\
  Para \(y(0) = x(-2)\), precisou do passado. \textbf{Com Memória}.
\item
  \textbf{Difícil:} \(y(t) = x(-t)\).\\
  Vamos testar \(t = 4\):
  \[y(4) = x(-4) \quad \text{(Precisou do passado)}\]
  Vamos testar \(t = -4\):
  \[y(-4) = x(4) \quad \text{(Precisou do futuro)}\]
  Tem \textbf{Memória}.
\end{itemize}

\subsection{4. Causalidade}\label{causalidade}

Causalidade significa que o sistema não pode prever o futuro.

\textbf{Algoritmo de Resolução:}

\begin{enumerate}
\def\labelenumi{\arabic{enumi}.}
\tightlist
\item
  \textbf{Passo 1:} Considere a saída num instante de tempo presente \(t_0\).
\item
  \textbf{Passo 2:} Verifique o argumento da entrada \(x(\cdot)\) na equação.
\item
  \textbf{Passo 3:} Se para qualquer valor escolhido de \(t_0\), o argumento for \textbf{maior} que \(t_0\), o sistema está olhando para o futuro. Logo, é \textbf{NÃO-CAUSAL}. Se o argumento for sempre \(\le t_0\), ele só olha para presente/passado e é \textbf{CAUSAL}.
\end{enumerate}

\begin{quote}
\textbf{Dica de Ouro / Macete Visual:}\\
Operadores \(+t\) dentro do `x', ou ``compressões/expansões'' temporais (\(x(2t)\)) costumam ser não-causais. Todo sistema sem memória é causal.
\end{quote}

\textbf{Exemplos Práticos:}

\begin{itemize}
\item
  \textbf{Fácil (Trivial):} \(y[n] = x[n] + x[n-1]\). Usa presente e passado. \textbf{Causal}.
\item
  \textbf{Médio:} \(y(t) = x(t+2)\). Preciso do futuro hoje. \textbf{Não-Causal}.
\item
  \textbf{Nível Prova/Lista (A Pegadinha do Oppenheim):} \(y(t) = t \cdot x(t)\).\\
  \emph{(Origem: Questão 5 da sua 1ª Avaliação)}\\
  Você avalia um instante \(t\): o valor da saída depende EXCLUSIVAMENTE da entrada \(x\) no mesmo exato momento \(t\). O coeficiente ``\(t\)'' multiplicando de fora é apenas um ganho variável, não muda o instante de leitura do sinal. Como não olha para o futuro, é \textbf{Causal} (e sem memória).
\end{itemize}

\subsection{5. LINEAR ou NÃO-LINEAR (Superposição)}\label{linear-ou-nuxe3o-linear-superposiuxe7uxe3o}

A superposição exige aditividade e homogeneidade.

\textbf{Algoritmo de Resolução:}

\begin{enumerate}
\def\labelenumi{\arabic{enumi}.}
\tightlist
\item
  \textbf{Passo 1:} Escreva as saídas individuais para duas entradas: \(y_1(t) = T\{x_1(t)\}\) e \(y_2(t) = T\{x_2(t)\}\).
\item
  \textbf{Passo 2:} Crie uma entrada combinada linearmente: \(x_3(t) = a \cdot x_1(t) + b \cdot x_2(t)\).
\item
  \textbf{Passo 3:} Calcule qual seria a saída do sistema \textbf{se} a entrada fosse \(x_3(t)\). Troque todo \(x(t)\) pela expressão inteira de \(x_3(t)\). Chame de \(y_3(t)\).
\item
  \textbf{Passo 4:} Calcule a soma linear das saídas individuais: \(y_{comb}(t) = a \cdot y_1(t) + b \cdot y_2(t)\).
\item
  \textbf{Passo 5:} Compare \(y_3(t)\) com \(y_{comb}(t)\). Se \(y_3(t) == y_{comb}(t)\), o sistema é \textbf{LINEAR}. Se forem diferentes, é \textbf{NÃO-LINEAR}.
\end{enumerate}

\begin{quote}
\textbf{Dica de Ouro / Macete Visual:}\\
Constantes somadas isoladas (ex: \(+3\)) ou o sinal de entrada enjaulado em funções não-lineares (\(x^2\), \(\sin(x)\)) destroem a linearidade.
\end{quote}

\textbf{Exemplos Práticos:}

\begin{itemize}
\item
  \textbf{Fácil (Trivial):} \(y(t) = x(t) + 3\).\\
  Entrada zero = \(x(t) = 0 \implies y(t) = 3\). Um sistema linear não pode produzir saída se não há entrada. Falhou na origem. \textbf{Não-Linear}.
\item
  \textbf{Nível Prova/Lista (Médio/Difícil):} O ganho variante \(y(t) = t \cdot x(t)\).\\
  \emph{(Origem: Questão 5 da sua 1ª Avaliação)}\\
  Combinação nas entradas:
  \[x_3(t) = a \cdot x_1(t) + b \cdot x_2(t)\]
  Passando pelo sistema (multiplicando a entrada inteira por t):
  \[y_3(t) = t \cdot [a \cdot x_1(t) + b \cdot x_2(t)]\]
  Distribuindo:
  \[y_3(t) = a \cdot [t \cdot x_1(t)] + b \cdot [t \cdot x_2(t)]\]
  \[y_3(t) = a \cdot y_1(t) + b \cdot y_2(t)\]
  Como a saída resultante separou direitinho e não quebrou a superposição, é \textbf{Linear}.
\end{itemize}

\subsection{6. Invariância no Tempo (LTI) ou Variante no Tempo}\label{invariuxe2ncia-no-tempo-lti-ou-variante-no-tempo}

Invariância no tempo significa que um atraso na entrada causa exatamente o mesmo atraso na saída.

\textbf{Algoritmo de Resolução:}

\begin{enumerate}
\def\labelenumi{\arabic{enumi}.}
\tightlist
\item
  \textbf{Passo 1:} Aplique uma entrada \(x(t)\) para obter a saída original \(y(t)\).
\item
  \textbf{Passo 2:} Crie um sinal de entrada atrasado artificialmente: \(x_1(t) = x(t - t_0)\).
\item
  \textbf{Passo 3:} Passe esse sinal \(x_1(t)\) pelo sistema (na equação, troque a notação \(x(\cdot)\) por \(x_1(\cdot)\)) e chame de \(y_1(t)\).
\item
  \textbf{Passo 4:} Pegue a expressão da saída original \(y(t)\) e substitua a variável \textbf{global} \(t\) por \((t - t_0)\) em \textbf{TODOS} os lugares onde \(t\) aparecer. Chame de \(y_2(t)\).
\item
  \textbf{Passo 5:} Compare \(y_1(t)\) com \(y_2(t)\). Se \(y_1(t) == y_2(t)\), o sistema é \textbf{INVARIANTE}. Caso contrário, é \textbf{VARIANTE}.
\end{enumerate}

\begin{quote}
\textbf{Dica de Ouro / Macete Visual:}\\
Variáveis de tempo (\(t\)) ``soltas'' do lado de fora do sinal indicam Variância no tempo. Alterações internas de relógio como \(x(2t)\) também indicam Variância.
\end{quote}

\textbf{Exemplos Práticos:}

\begin{itemize}
\item
  \textbf{Fácil (Trivial):} \(y(t) = 2x(t)\).\\
  Atrasando a entrada e a saída, ambas dão \(2x(t-t_0)\). \textbf{Invariante}.
\item
  \textbf{Nível Prova/Lista (A clássica de Variância):} \(y(t) = t \cdot x(t)\).\\
  \emph{(Origem: Questão 5 da sua 1ª Avaliação)}\\
  Passo 1 (atrasa SÓ a entrada \(x\)):
  \[y_1(t) = t \cdot x(t-t_0)\]
  Passo 2 (atrasa a variável de tempo `t' global do sistema):
  \[y_2(t) = (t-t_0) \cdot x(t-t_0)\]
  Avaliando: \(y_1(t) \neq y_2(t)\). Logo, \textbf{Variante no Tempo}.
\item
  \textbf{Difícil (Compressão no Tempo):} \(y(t) = x(2t)\).\\
  Passo 1 (Atraso na entrada):
  \[y_1(t) = x(2t - t_0)\]
  Passo 2 (trocar a variável de tempo global \(t\) por \(t-t_0\)):
  \[y_2(t) = x(2(t-t_0))\]
  \[y_2(t) = x(2t - 2t_0)\]
  Comparando:
  \[x(2t - t_0) \neq x(2t - 2t_0)\]
  Logo, é \textbf{Variante no tempo}.
\end{itemize}

\subsection{7. Estabilidade BIBO}\label{estabilidade-bibo}

Se entra um valor limitado (finito), tem que sair um valor limitado (finito).

\textbf{Algoritmo de Resolução:}

\begin{enumerate}
\def\labelenumi{\arabic{enumi}.}
\tightlist
\item
  \textbf{Passo 1:} Assuma, por hipótese, que a entrada é limitada. Escreva: \(|x(t)| \le B_x < \infty\) para todo \(t\).
\item
  \textbf{Passo 2:} Aplique o módulo na equação inteira da saída: \(|y(t)| = |T\{x(t)\}|\).
\item
  \textbf{Passo 3:} Use propriedades de módulo e desigualdade triangular para isolar o termo \(|x(t)|\).
\item
  \textbf{Passo 4:} Substitua \(|x(t)|\) por \(B_x\).
\item
  \textbf{Passo 5:} Analise a expressão. Se houver alguma entrada limitada (mesmo que seja um caso apenas) que faça a saída ir para infinito (ex: divisão por zero ou integração de constante infinita), é \textbf{INSTÁVEL}. Se o limite for sempre finito, é \textbf{ESTÁVEL}.
\end{enumerate}

\begin{quote}
\textbf{Macete Visual:}\\
Integradores com limites abertos (\(-\infty\) até \(t\)) tendem a instabilizar com entradas constantes. Funções divididas por tempo (\(1/t\)) podem instabilizar em zero.
\end{quote}

\textbf{Exemplos Práticos:}

\begin{itemize}
\item
  \textbf{Fácil (Trivial):} Divisor \(y(t) = \frac{1}{2} x(t)\).
  \[|y(t)| = \frac{1}{2}|x(t)| \le \frac{1}{2} B_x < \infty\]. \textbf{Estável}.
\item
  \textbf{Nível Prova/Lista (Pegadinha do ``t''):} O ganho variante \(y(t) = t \cdot x(t)\).\\
  \emph{(Origem: Questão 5 da sua 1ª Avaliação)}\\
  Assuma que aplicamos um degrau unitário \(x(t) = 1\) (sinal super comportado e limitado, \(B_x=1\)).\\
  A saída será:
  \[y(t) = t \cdot (1) = t\]
  Se avaliarmos quando \(t \to \infty\), a saída \(y(t) \to \infty\).\\
  Ou seja, entrou um sinal com limite e a saída estourou pro infinito por conta do \(t\) livre! Sistema \textbf{Instável}.
\end{itemize}
